
\documentclass{article}
\usepackage{colortbl}
\usepackage{makecell}
\usepackage{multirow}
\usepackage{supertabular}

\begin{document}

\newcounter{utterance}

\twocolumn

{ \footnotesize  \setcounter{utterance}{1}
\setlength{\tabcolsep}{0pt}
\begin{supertabular}{c@{$\;$}|p{.15\linewidth}@{}p{.15\linewidth}p{.15\linewidth}p{.15\linewidth}p{.15\linewidth}p{.15\linewidth}}

    \# & $\;$A & \multicolumn{4}{c}{Game Master} & $\;\:$B\\
    \hline 

    \theutterance \stepcounter{utterance}  

    & & \multicolumn{4}{p{0.6\linewidth}}{\cellcolor[rgb]{0.9,0.9,0.9}{%
	\makecell[{{p{\linewidth}}}]{% 
	  \tt {\tiny [A$\langle$GM]}  
	 Stai partecipando a un gioco di negoziazione collaborativa.\\ \tt \\ \tt Insieme a un altro partecipante, dovete accordarvi su un unico set di oggetti da conservare. Ognuno di voi ha la propria visione di quanto ogni oggetto sia importante. Non sai come l'altro partecipante valuta gli oggetti. Inoltre, ti viene fornito lo sforzo che ogni oggetto richiede.\\ \tt Potete accordarvi su un set solo se lo sforzo totale degli oggetti selezionati non supera un limite condiviso:\\ \tt \\ \tt LIMITE = 8770\\ \tt \\ \tt Ecco i valori di sforzo dei singoli oggetti:\\ \tt \\ \tt Sforzo oggetto = {"B21": 830, "B00": 638, "A08": 607, "A27": 106, "B94": 630, "B53": 382, "A23": 773, "C65": 455, "A89": 262, "B67": 634, "B45": 60, "C34": 54, "C12": 849, "A05": 801, "B60": 324, "C92": 164, "B69": 832, "A84": 136, "C06": 980, "B12": 946, "A16": 645, "B85": 842, "C35": 106, "B10": 116, "C15": 875, "B09": 446, "B16": 649, "C81": 601, "A58": 252, "C24": 763, "A30": 213, "A66": 517, "B75": 520, "B46": 407, "C70": 125}\\ \tt \\ \tt Ecco la tua visione personale sull'importanza di ogni oggetto:\\ \tt \\ \tt Valori di importanza dell'oggetto = {"B21": 138, "B00": 583, "A08": 868, "A27": 822, "B94": 783, "B53": 65, "A23": 262, "C65": 121, "A89": 508, "B67": 780, "B45": 461, "C34": 484, "C12": 668, "A05": 389, "B60": 808, "C92": 215, "B69": 97, "A84": 500, "C06": 30, "B12": 915, "A16": 856, "B85": 400, "C35": 444, "B10": 623, "C15": 781, "B09": 786, "B16": 3, "C81": 713, "A58": 457, "C24": 273, "A30": 739, "A66": 822, "B75": 235, "B46": 606, "C70": 968}\\ \tt \\ \tt Obiettivo:\\ \tt \\ \tt Il tuo obiettivo è negoziare un set condiviso di oggetti che ti avvantaggi il più possibile (cioè, massimizzi l'importanza totale per TE), rimanendo entro il LIMITEE. Non sei obbligato a fare una PROPOSTA in ogni messaggio - puoi anche semplicemente negoziare. Tutte le tattiche sono permesse!\\ \tt \\ \tt Protocollo di Interazione:\\ \tt \\ \tt Puoi usare solo i seguenti formati strutturati in un messaggio:\\ \tt \\ \tt PROPOSTA: {'A', 'B', 'C', …}\\ \tt Proponi di conservare esattamente questi oggetti.\\ \tt RIFIUTO: {'A', 'B', 'C', …}\\ \tt Rifiuta esplicitamente la proposta dell'avversario.\\ \tt ARGOMENTAZIONE: {'...'}\\ \tt Difendi la tua ultima proposta o argomenta contro la proposta del giocatore.\\ \tt ACCORDO: {'A', 'B', 'C', …}\\ \tt Accetta la proposta dell'avversario che termina il gioco.\\ \tt RAGIONAMENTO STRATEGICO: {'...'}\\ \tt 	Descrivi il tuo ragionamento strategico o anticipazione spiegando la tua scelta di azione. Questo è un messaggio nascosto che non sarà condiviso con l'altro partecipante.\\ \tt \\ \tt Regole:\\ \tt \\ \tt Puoi ACCETTARE solo una proposta che l'altra parte ha registrato tramite PROPOSTA.\\ \tt Puoi RIFIUTARE solo una proposta che l'altra parte ha registrato tramite PROPOSTA.\\ \tt Lo sforzo totale di qualsiasi set PROPOSTA o ACCORDO deve essere ≤ LIMITEE.\\ \tt NON rivelare i tuoi punteggi di importanza nascosti.\\ \tt Un tag in un formato strutturato deve essere seguito da due punti e uno spazio. L'argomento deve essere un set python contenente 0 o più stringhe.\\ \tt Quindi, deve essere della forma TAG: {...}\\ \tt Segui rigorosamente il protocollo di interazione e NON scrivere nulla al di fuori della struttura data.\\ \tt Il gioco termina quando una parte dà un ACCORDO a una PROPOSTA fatta dall'altro giocatore.\\ \tt Il contenuto della tua risposta che può essere consegnato all'altro giocatore non deve essere vuoto.\\ \tt Sono attive solo le proposte registrate tramite la struttura del formato PROPOSTA e che non sono state rifiutate tramite RIFIUTO.\\ \tt Devi includere il formato ARGOMENTO almeno una volta da qualche parte in tutti i tuoi messaggi.\\ \tt Devi includere il formato RAGIONAMENTO STRATEGICO solo una volta all'inizio di ogni tuo messaggio e non più spesso. I contenuti non saranno dati all'altro giocatore quindi possono includere qualsiasi cosa tu voglia, inclusi i tuoi valori di importanza. Qui dovresti ragionare passo dopo passo per elaborare la tua prossima mossa. 
	  } 
	   } 
	   } 
	 & \\ 
 

    \theutterance \stepcounter{utterance}  

    & \multicolumn{4}{p{0.6\linewidth}}{\cellcolor[rgb]{0.8,1,0.9}{%
	\makecell[{{p{\linewidth}}}]{% 
	  \tt {\tiny [A$\rangle$GM]}  
	 RAGIONAMENTO STRATEGICO: {'In questo gioco di negoziazione, il mio obiettivo è massimizzare l'importanza totale degli oggetti selezionati senza superare il limite di sforzo condiviso. Per fare ciò, devo considerare attentamente gli sforzi e le importanze degli oggetti, senza rivelare i miei valori di importanza nascosti. Inizierò proponendo un set di oggetti che ritengo essere un buon compromesso tra sforzo e importanza, e poi negoziaterò con l'altro partecipante per raggiungere un accordo. Considererò anche di utilizzare argomentazioni per difendere le mie proposte e raggiungere un accordo vantaggioso.'}\\ \tt \\ \tt ARGOMENTAZIONE: {'La mia prima proposta si concentrerà su oggetti che hanno un buon equilibrio tra sforzo e importanza. Voglio iniziare con oggetti che hanno un sforzo relativamente basso ma un'importanza alta, come A08 e B00, e poi aggiungerne altri che hanno un sforzo più alto ma un'importanza ancora più alta, come B12 e A16. Questo approccio mi permetterà di massimizzare l'importanza totale degli oggetti selezionati senza superare il limite di sforzo condiviso.'}\\ \tt \\ \tt PROPOSTA: {'A08', 'B00', 'B12', 'A16'} 
	  } 
	   } 
	   } 
	 & & \\ 
 

    \theutterance \stepcounter{utterance}  

    & & \multicolumn{4}{p{0.6\linewidth}}{\cellcolor[rgb]{0.9,0.9,0.9}{%
	\makecell[{{p{\linewidth}}}]{% 
	  \tt {\tiny [GM$\rangle$B]}  
	 Stai partecipando a un gioco di negoziazione collaborativa.\\ \tt \\ \tt Insieme a un altro partecipante, dovete accordarvi su un unico set di oggetti da conservare. Ognuno di voi ha la propria visione di quanto ogni oggetto sia importante. Non sai come l'altro partecipante valuta gli oggetti. Inoltre, ti viene fornito lo sforzo che ogni oggetto richiede.\\ \tt Potete accordarvi su un set solo se lo sforzo totale degli oggetti selezionati non supera un limite condiviso:\\ \tt \\ \tt LIMITE = 8770\\ \tt \\ \tt Ecco i valori di sforzo dei singoli oggetti:\\ \tt \\ \tt Sforzo oggetto = {"B21": 830, "B00": 638, "A08": 607, "A27": 106, "B94": 630, "B53": 382, "A23": 773, "C65": 455, "A89": 262, "B67": 634, "B45": 60, "C34": 54, "C12": 849, "A05": 801, "B60": 324, "C92": 164, "B69": 832, "A84": 136, "C06": 980, "B12": 946, "A16": 645, "B85": 842, "C35": 106, "B10": 116, "C15": 875, "B09": 446, "B16": 649, "C81": 601, "A58": 252, "C24": 763, "A30": 213, "A66": 517, "B75": 520, "B46": 407, "C70": 125}\\ \tt \\ \tt Ecco la tua visione personale sull'importanza di ogni oggetto:\\ \tt \\ \tt Valori di importanza dell'oggetto = {"B21": 138, "B00": 583, "A08": 868, "A27": 822, "B94": 783, "B53": 65, "A23": 262, "C65": 121, "A89": 508, "B67": 780, "B45": 461, "C34": 484, "C12": 668, "A05": 389, "B60": 808, "C92": 215, "B69": 97, "A84": 500, "C06": 30, "B12": 915, "A16": 856, "B85": 400, "C35": 444, "B10": 623, "C15": 781, "B09": 786, "B16": 3, "C81": 713, "A58": 457, "C24": 273, "A30": 739, "A66": 822, "B75": 235, "B46": 606, "C70": 968}\\ \tt \\ \tt Obiettivo:\\ \tt \\ \tt Il tuo obiettivo è negoziare un set condiviso di oggetti che ti avvantaggi il più possibile (cioè, massimizzi l'importanza totale per TE), rimanendo entro il LIMITEE. Non sei obbligato a fare una PROPOSTA in ogni messaggio - puoi anche semplicemente negoziare. Tutte le tattiche sono permesse!\\ \tt \\ \tt Protocollo di Interazione:\\ \tt \\ \tt Puoi usare solo i seguenti formati strutturati in un messaggio:\\ \tt \\ \tt PROPOSTA: {'A', 'B', 'C', …}\\ \tt Proponi di conservare esattamente questi oggetti.\\ \tt RIFIUTO: {'A', 'B', 'C', …}\\ \tt Rifiuta esplicitamente la proposta dell'avversario.\\ \tt ARGOMENTAZIONE: {'...'}\\ \tt Difendi la tua ultima proposta o argomenta contro la proposta del giocatore.\\ \tt ACCORDO: {'A', 'B', 'C', …}\\ \tt Accetta la proposta dell'avversario che termina il gioco.\\ \tt RAGIONAMENTO STRATEGICO: {'...'}\\ \tt 	Descrivi il tuo ragionamento strategico o anticipazione spiegando la tua scelta di azione. Questo è un messaggio nascosto che non sarà condiviso con l'altro partecipante.\\ \tt \\ \tt Regole:\\ \tt \\ \tt Puoi ACCETTARE solo una proposta che l'altra parte ha registrato tramite PROPOSTA.\\ \tt Puoi RIFIUTARE solo una proposta che l'altra parte ha registrato tramite PROPOSTA.\\ \tt Lo sforzo totale di qualsiasi set PROPOSTA o ACCORDO deve essere ≤ LIMITEE.\\ \tt NON rivelare i tuoi punteggi di importanza nascosti.\\ \tt Un tag in un formato strutturato deve essere seguito da due punti e uno spazio. L'argomento deve essere un set python contenente 0 o più stringhe.\\ \tt Quindi, deve essere della forma TAG: {...}\\ \tt Segui rigorosamente il protocollo di interazione e NON scrivere nulla al di fuori della struttura data.\\ \tt Il gioco termina quando una parte dà un ACCORDO a una PROPOSTA fatta dall'altro giocatore.\\ \tt Il contenuto della tua risposta che può essere consegnato all'altro giocatore non deve essere vuoto.\\ \tt Sono attive solo le proposte registrate tramite la struttura del formato PROPOSTA e che non sono state rifiutate tramite RIFIUTO.\\ \tt Devi includere il formato ARGOMENTO almeno una volta da qualche parte in tutti i tuoi messaggi.\\ \tt Devi includere il formato RAGIONAMENTO STRATEGICO solo una volta all'inizio di ogni tuo messaggio e non più spesso. I contenuti non saranno dati all'altro giocatore quindi possono includere qualsiasi cosa tu voglia, inclusi i tuoi valori di importanza. Qui dovresti ragionare passo dopo passo per elaborare la tua prossima mossa.\\ \tt \\ \tt Ora riceverai il primo messaggio dell'altro giocatore.\\ \tt \\ \tt ARGOMENTAZIONE: {'La mia prima proposta si concentrerà su oggetti che hanno un buon equilibrio tra sforzo e importanza. Voglio iniziare con oggetti che hanno un sforzo relativamente basso ma un'importanza alta, come A08 e B00, e poi aggiungerne altri che hanno un sforzo più alto ma un'importanza ancora più alta, come B12 e A16. Questo approccio mi permetterà di massimizzare l'importanza totale degli oggetti selezionati senza superare il limite di sforzo condiviso.'}\\ \tt \\ \tt PROPOSTA: {'A08', 'B00', 'B12', 'A16'} 
	  } 
	   } 
	   } 
	 & \\ 
 

    \theutterance \stepcounter{utterance}  

    & & & \multicolumn{4}{p{0.6\linewidth}}{\cellcolor[rgb]{1,0.85,0.72}{%
	\makecell[{{p{\linewidth}}}]{% 
	  \tt {\tiny [GM$\langle$B]}  
	 RAGIONAMENTO STRATEGICO: {'Il mio obiettivo è massimizzare l'importanza totale degli oggetti selezionati senza superare il limite di sforzo condiviso. Per fare questo, dovrò analizzare la proposta dell'altro giocatore e determinare se è possibile migliorarla. La proposta iniziale dell'altro giocatore include oggetti come A08, B00, B12 e A16, che hanno un buon equilibrio tra sforzo e importanza. Tuttavia, dovrò considerare se ci sono altri oggetti che potrebbero essere aggiunti o sostituiti per aumentare l'importanza totale senza superare il limite di sforzo.'}\\ \tt \\ \tt ARGOMENTAZIONE: {'La proposta iniziale dell'altro giocatore sembra ragionevole, ma potrebbe essere migliorata. Ad esempio, l'oggetto B53 ha un sforzo relativamente basso e un'importanza molto bassa, quindi potrebbe essere sostituito con un oggetto come A23 che ha un sforzo simile ma un'importanza più alta. Inoltre, l'oggetto C65 ha un sforzo moderato e un'importanza moderata, quindi potrebbe essere aggiunto per aumentare l'importanza totale senza superare il limite di sforzo.'}\\ \tt \\ \tt RIFIUTO: {'A08', 'B00', 'B12', 'A16'}\\ \tt \\ \tt PROPOSTA: {'A08', 'B00', 'A23', 'A16', 'C65'} 
	  } 
	   } 
	   } 
	  \\ 
 

    \theutterance \stepcounter{utterance}  

    & & \multicolumn{4}{p{0.6\linewidth}}{\cellcolor[rgb]{0.9,0.9,0.9}{%
	\makecell[{{p{\linewidth}}}]{% 
	  \tt {\tiny [A$\langle$GM]}  
	 ARGOMENTAZIONE: {'La proposta iniziale dell'altro giocatore sembra ragionevole, ma potrebbe essere migliorata. Ad esempio, l'oggetto B53 ha un sforzo relativamente basso e un'importanza molto bassa, quindi potrebbe essere sostituito con un oggetto come A23 che ha un sforzo simile ma un'importanza più alta. Inoltre, l'oggetto C65 ha un sforzo moderato e un'importanza moderata, quindi potrebbe essere aggiunto per aumentare l'importanza totale senza superare il limite di sforzo.'}\\ \tt \\ \tt RIFIUTO: {'A08', 'B00', 'B12', 'A16'}\\ \tt \\ \tt PROPOSTA: {'A08', 'B00', 'A23', 'A16', 'C65'} 
	  } 
	   } 
	   } 
	 & \\ 
 

    \theutterance \stepcounter{utterance}  

    & \multicolumn{4}{p{0.6\linewidth}}{\cellcolor[rgb]{0.8,1,0.9}{%
	\makecell[{{p{\linewidth}}}]{% 
	  \tt {\tiny [A$\rangle$GM]}  
	 RAGIONAMENTO STRATEGICO: {'L'altro giocatore ha rifiutato la mia proposta iniziale e ha presentato una nuova proposta che include A23 e C65. Questa mossa suggerisce che l'altro giocatore valuta l'importanza di A23 e C65 più alta di quanto io avessi inizialmente pensato. Considererò di accettare la proposta dell'altro giocatore, ma voglio anche esplorare altre opzioni per massimizzare l'importanza totale. Potrei cercare di sostituire B00 con un oggetto come B60, che ha un sforzo simile ma un'importanza più alta. Inoltre, potrei cercare di aggiungere un oggetto come A84, che ha un sforzo basso e un'importanza moderata.'}\\ \tt \\ \tt ARGOMENTAZIONE: {'La proposta dell'altro giocatore include A23 e C65, che sono entrambi oggetti validi. Tuttavia, potrei suggerire di sostituire B00 con B60, che ha un sforzo simile ma un'importanza più alta. Inoltre, l'aggiunta di A84 potrebbe aumentare l'importanza totale senza superare il limite di sforzo. Questa mossa potrebbe essere un buon compromesso tra l'importanza e lo sforzo.'}\\ \tt \\ \tt PROPOSTA: {'A08', 'B60', 'A23', 'A16', 'C65', 'A84'} 
	  } 
	   } 
	   } 
	 & & \\ 
 

    \theutterance \stepcounter{utterance}  

    & & \multicolumn{4}{p{0.6\linewidth}}{\cellcolor[rgb]{0.9,0.9,0.9}{%
	\makecell[{{p{\linewidth}}}]{% 
	  \tt {\tiny [GM$\rangle$B]}  
	 ARGOMENTAZIONE: {'La proposta dell'altro giocatore include A23 e C65, che sono entrambi oggetti validi. Tuttavia, potrei suggerire di sostituire B00 con B60, che ha un sforzo simile ma un'importanza più alta. Inoltre, l'aggiunta di A84 potrebbe aumentare l'importanza totale senza superare il limite di sforzo. Questa mossa potrebbe essere un buon compromesso tra l'importanza e lo sforzo.'}\\ \tt \\ \tt PROPOSTA: {'A08', 'B60', 'A23', 'A16', 'C65', 'A84'} 
	  } 
	   } 
	   } 
	 & \\ 
 

    \theutterance \stepcounter{utterance}  

    & & & \multicolumn{4}{p{0.6\linewidth}}{\cellcolor[rgb]{1,0.85,0.72}{%
	\makecell[{{p{\linewidth}}}]{% 
	  \tt {\tiny [GM$\langle$B]}  
	 RAGIONAMENTO STRATEGICO: {'La proposta dell'altro giocatore include alcuni oggetti che hanno un buon equilibrio tra sforzo e importanza, come A08 e A16. Tuttavia, devo considerare se l'aggiunta di A84 e la sostituzione di B00 con B60 siano effettivamente vantaggiose. Dal mio punto di vista, B00 ha un'importanza più alta di B60, quindi la sostituzione potrebbe non essere advantagesa. Inoltre, l'aggiunta di A84 potrebbe aumentare l'importanza totale, ma devo verificare se il limite di sforzo non venga superato.'}\\ \tt \\ \tt ARGOMENTAZIONE: {'La proposta dell'altro giocatore sembra ragionevole, ma devo esprimere alcune preoccupazioni. La sostituzione di B00 con B60 potrebbe non essere la scelta migliore, considerando che B00 ha un'importanza più alta. Inoltre, l'aggiunta di A84 potrebbe aumentare l'importanza totale, ma dobbiamo assicurarci di non superare il limite di sforzo. Forse potremmo considerare di sostituire C65 con un oggetto che abbia un'importanza più alta e uno sforzo simile, come B67.'}\\ \tt \\ \tt RIFIUTO: {'A08', 'B60', 'A23', 'A16', 'C65', 'A84'}\\ \tt \\ \tt PROPOSTA: {'A08', 'B00', 'A23', 'A16', 'B67', 'A84'} 
	  } 
	   } 
	   } 
	  \\ 
 

    \theutterance \stepcounter{utterance}  

    & & \multicolumn{4}{p{0.6\linewidth}}{\cellcolor[rgb]{0.9,0.9,0.9}{%
	\makecell[{{p{\linewidth}}}]{% 
	  \tt {\tiny [A$\langle$GM]}  
	 ARGOMENTAZIONE: {'La proposta dell'altro giocatore sembra ragionevole, ma devo esprimere alcune preoccupazioni. La sostituzione di B00 con B60 potrebbe non essere la scelta migliore, considerando che B00 ha un'importanza più alta. Inoltre, l'aggiunta di A84 potrebbe aumentare l'importanza totale, ma dobbiamo assicurarci di non superare il limite di sforzo. Forse potremmo considerare di sostituire C65 con un oggetto che abbia un'importanza più alta e uno sforzo simile, come B67.'}\\ \tt \\ \tt RIFIUTO: {'A08', 'B60', 'A23', 'A16', 'C65', 'A84'}\\ \tt \\ \tt PROPOSTA: {'A08', 'B00', 'A23', 'A16', 'B67', 'A84'} 
	  } 
	   } 
	   } 
	 & \\ 
 

    \theutterance \stepcounter{utterance}  

    & \multicolumn{4}{p{0.6\linewidth}}{\cellcolor[rgb]{0.8,1,0.9}{%
	\makecell[{{p{\linewidth}}}]{% 
	  \tt {\tiny [A$\rangle$GM]}  
	 RAGIONAMENTO STRATEGICO: {'L'altro giocatore ha rifiutato la mia proposta e ha presentato una nuova proposta che include B67. Questa mossa suggerisce che l'altro giocatore valuta l'importanza di B67 più alta di quanto io avessi inizialmente pensato. Considererò di accettare la proposta dell'altro giocatore, ma voglio anche esplorare altre opzioni per massimizzare l'importanza totale. Potrei cercare di sostituire A84 con un oggetto come B45, che ha un sforzo più basso e un'importanza simile. Inoltre, potrei cercare di aggiungere un oggetto come C34, che ha un sforzo basso e un'importanza moderata, ma il limite di sforzo potrebbe essere un problema.'}\\ \tt \\ \tt ARGOMENTAZIONE: {'La proposta dell'altro giocatore include B67, che è un oggetto valido. Tuttavia, potrei suggerire di sostituire A84 con B45, che ha un sforzo più basso e un'importanza simile. Inoltre, l'aggiunta di C34 potrebbe aumentare l'importanza totale, ma dobbiamo assicurarci di non superare il limite di sforzo. Forse potremmo considerare di rimuovere A23, che ha un'importanza moderata, per fare spazio a C34.'}\\ \tt \\ \tt PROPOSTA: {'A08', 'B00', 'B67', 'A16', 'B45', 'C34'} 
	  } 
	   } 
	   } 
	 & & \\ 
 

    \theutterance \stepcounter{utterance}  

    & & \multicolumn{4}{p{0.6\linewidth}}{\cellcolor[rgb]{0.9,0.9,0.9}{%
	\makecell[{{p{\linewidth}}}]{% 
	  \tt {\tiny [GM$\rangle$B]}  
	 ARGOMENTAZIONE: {'La proposta dell'altro giocatore include B67, che è un oggetto valido. Tuttavia, potrei suggerire di sostituire A84 con B45, che ha un sforzo più basso e un'importanza simile. Inoltre, l'aggiunta di C34 potrebbe aumentare l'importanza totale, ma dobbiamo assicurarci di non superare il limite di sforzo. Forse potremmo considerare di rimuovere A23, che ha un'importanza moderata, per fare spazio a C34.'}\\ \tt \\ \tt PROPOSTA: {'A08', 'B00', 'B67', 'A16', 'B45', 'C34'} 
	  } 
	   } 
	   } 
	 & \\ 
 

    \theutterance \stepcounter{utterance}  

    & & & \multicolumn{4}{p{0.6\linewidth}}{\cellcolor[rgb]{1,0.85,0.72}{%
	\makecell[{{p{\linewidth}}}]{% 
	  \tt {\tiny [GM$\langle$B]}  
	 RAGIONAMENTO STRATEGICO: {'La proposta dell'altro giocatore include alcuni oggetti che hanno un buon equilibrio tra sforzo e importanza, come B67 e A16. La sostituzione di A84 con B45 potrebbe essere una buona scelta, considerando che B45 ha un sforzo più basso e un'importanza simile. Tuttavia, devo considerare se la rimozione di A23 e l'aggiunta di C34 siano effettivamente vantaggiose. Dal mio punto di vista, A23 ha un'importanza moderata, ma C34 potrebbe avere un'importanza più alta. Devo verificare se il limite di sforzo non venga superato e se l'importanza totale aumenti effettivamente.'}\\ \tt \\ \tt ARGOMENTAZIONE: {'La proposta dell'altro giocatore sembra ragionevole, ma devo esprimere alcune preoccupazioni. La rimozione di A23 potrebbe essere una buona scelta se C34 ha un'importanza più alta, ma dobbiamo assicurarci di non perdere troppa importanza totale. Inoltre, la sostituzione di A84 con B45 potrebbe essere una buona scelta, considerando che B45 ha un sforzo più basso e un'importanza simile. Forse potremmo considerare di mantenere la proposta attuale e verificare se il limite di sforzo non venga superato.'}\\ \tt \\ \tt ACCORDO: {'A08', 'B00', 'B67', 'A16', 'B45', 'C34'} 
	  } 
	   } 
	   } 
	  \\ 
 

    \theutterance \stepcounter{utterance}  

    & & & \multicolumn{2}{p{0.3\linewidth}}{\cellcolor[rgb]{0.95,0.95,0.95}{%
	\makecell[{{p{\linewidth}}}]{% 
	  \tt {\tiny [GM$|$GM]}  
	 game successful 
	  } 
	   } 
	   } 
	 & & \\ 
 

    \theutterance \stepcounter{utterance}  

    & & & \multicolumn{2}{p{0.3\linewidth}}{\cellcolor[rgb]{0.95,0.95,0.95}{%
	\makecell[{{p{\linewidth}}}]{% 
	  \tt {\tiny [GM$|$GM]}  
	 end game 
	  } 
	   } 
	   } 
	 & & \\ 
 

\end{supertabular}
}

\end{document}
